\chapter{项目管理}

排期时，我们先问每轮要交什么、由谁接手、什么时候可以一起联调。这几件事对应课程讲的 4P。产品的范围写进用户故事，人员分工落实到工作包，过程约定接口和验收的顺序，项目层面再跟踪工时、交付日期和风险。跨端接口、交易或外部服务有问题时，我们按它对当轮目标的影响安排修复。

\section{任务计划安排}

\subsection{工作分解（WBS）}

把“顾客能下单”分给一个人，到了联调仍说不清哪一步算做完。我们按 WBS 的“100\% 规则”把需求、测试和交付文档都拆进任务，最小任务尽量安排在 4--8 小时左右。图~\ref{fig:mgmt-wbs-tree} 画出了三层结构。顺着用户故事能找到要完成的改动，手头的任务也能归回当轮交付目标。

\begin{reportpoints}
\item 第一层按需求和工程基础分组，覆盖 \code{US001}--\code{US022} 的 22 个用户故事，以及环境、测试和交付文档。
\item 第二层按 Sprint 列出交付模块，任务编号为 \code{T101}--\code{T308}。接口、数据库、前端与质量工作依实际联调顺序建立依赖。
\item 第三层写明可单独验收的改动，如一项 Controller 接口、一个 SQL 迁移、一个 Vue 组件或一组回归，记录负责人、估算工时和完成标准。
\end{reportpoints}

\begin{figure}[!htbp]
  \centering
  \resizebox{.94\linewidth}{!}{%
\begin{tikzpicture}[
    node distance=0.7cm and 0.45cm,
    root/.style={rectangle,rounded corners=4pt,draw=tjublue,very thick,fill=tjublue,text=white,font=\small\bfseries,minimum width=5.6cm,minimum height=0.9cm,align=center},
    l1/.style={rectangle,rounded corners=3pt,draw=tjublue,thick,fill=lightblue,font=\scriptsize\bfseries,minimum width=2.8cm,minimum height=0.9cm,align=center},
    l2/.style={rectangle,rounded corners=2pt,draw=mutedgray,fill=white,font=\scriptsize,minimum width=2.7cm,minimum height=1.55cm,align=left,inner sep=4pt},
    edge/.style={-{Stealth[length=1.8mm]},thick,tjublue!80}
  ]
    \node[root] (top) {轻量级外卖服务平台项目（230 h 基准 / 118 故事点）};
    \node[l1,below left=1.0cm and 2.5cm of top] (s1) {WBS-1.0 基础设施与\\Sprint 1 核心下单（80 h）};
    \node[l1,below=1.0cm of top] (s2) {WBS-2.0 履约、配送\\与资产闭环（68 h）};
    \node[l1,below right=1.0cm and 2.5cm of top] (s3) {WBS-3.0 管理、AI、地图\\与质量交付（82 h）};

    \draw[edge] (top.south) -- ++(0,-0.45) -| (s1.north);
    \draw[edge] (top.south) -- (s2.north);
    \draw[edge] (top.south) -- ++(0,-0.45) -| (s3.north);

    \node[l2,below=0.55cm of s1] (s1d) {1.1 SRS 与 17 张建模图\\1.2 \code{elm_v2.sql} 与迁移\\1.3 认证、店铺与菜品\\1.4 购物车与幂等下单\\1.5 顾客主链路与单测};
    \node[l2,below=0.55cm of s2] (s2d) {2.1 商家接单出餐管理\\2.2 骑手乐观锁抢单\\2.3 钱包积分券流水账\\2.4 评价、收藏与通知\\2.5 登录视觉与首次联调};
    \node[l2,below=0.55cm of s3] (s3d) {3.1 管理员审核与仲裁\\3.2 多模态 AI 与熔断\\3.3 高德地图与坐标距离\\3.4 四套主题与数据看板\\3.5 575 项回归与容器交付};

    \draw[edge] (s1) -- (s1d);
    \draw[edge] (s2) -- (s2d);
    \draw[edge] (s3) -- (s3d);
  \end{tikzpicture}
}
  \caption{项目三级工作分解结构（WBS）顶层树状图}\label{fig:mgmt-wbs-tree}
\end{figure}

以 \code{T105} 为例，顾客要能提交和取消订单，订单事务因此拆成三个能单独验的工作包。\code{T105-1} 检查优惠、配送费和自取计价，估算 4.0~h。\code{T105-2} 检查扣库存、写快照和防重提交，估算 4.5~h。\code{T105-3} 检查取消时的权限以及库存、资产恢复，估算 3.5~h。三项共 12.0~h，能与 Sprint 1 任务表对上。联调时先跑各自用例，再用一笔完整订单把它们接起来。

在 Sprint 2，\code{T202} 先完成商家接单、拒单和出餐，再接入骑手任务池与乐观锁抢单。两类状态需要在同一笔订单上衔接，所以商家前端和骑手工作台都依赖该模块。在 Sprint 3，\code{T307} 将页面返回回归与云端部署拆成独立任务，前者交付路由测试，后者交付脚本和部署文档。其余任务与依赖见附录~\ref{app:process}，三级拆解示例见表~\ref{tab:wbs-detail}。

\paragraph{跨模块任务的依赖}
前置任务决定哪一天可以联调。Sprint 1 先有数据库基线 \code{T102}，认证和店铺查询才能按同一份数据定义开发。顾客首页对接店铺查询 \code{T104}，购物车结算和订单列表则要等 \code{T105} 可以建单。页面可以先按约定开发，到接口可用时再核对完整操作。

同样的安排用于第三轮。管理员端看板要读取平台统计结果，先排权限与平台治理，再安排统计接口和导出。多模态交互的页面则同结构化助手服务对接，前端主题、动效和页面内容可以并行推进。\code{T308} 的全链路验收排在这几项之后，一次回归可以检查组合在一起的问题。若需要返工，也可沿这张任务表核对由谁负责接口、页面或测试数据。

图~\ref{fig:mgmt-wbs-tree} 的三条分支分别计划投入 80、68 和 82 小时，合计 230 小时。各阶段的工作量也包含测试和交付。首轮在下单功能之后安排单测，第二轮在履约、资产和评价等模块之后集中联调，末轮另外预留全量回归和容器交付的工时。这样核查某项任务时，能同时看到它所在的交付范围和它所依赖的工作。

\subsection{任务估算与偏差分析}

同一项任务，单独写代码和几端联调往往用时不同。我们使用 PERT 三点估算，由主责人与 Scrum Master 给出乐观工时 $O$、最可能工时 $M$ 和悲观工时 $P$，依式~(\ref{eq:pert}) 计算期望工时 $E$ 及标准差 $\sigma$。$O$、$P$ 的权重各为 1，$M$ 的权重为 4。
\begin{equation}\label{eq:pert}
  E = \frac{O + 4M + P}{6}, \qquad \sigma = \frac{P - O}{6}
\end{equation}
以认证、检索、购物车和下单的 \code{T103--105} 为例，表~\ref{tab:pert-estimation} 记录了 24、30、36~h 三种情况，期望 30~h、标准差 2~h，实际用时 31.5~h。这里多用的 1.5~h 是项目总差额 8.0~h 的一部分。三轮计划投入 230.0~h，实际投入 238.0~h，偏差率为 $+3.48\%$。

认证、检索和交易从乐观到悲观相差 12~h，只跑首轮认证和下单测试则相差 4~h。并行工作的几组人能不能拿到同一个可跑的版本，会影响前一组工时。实际超出估算时，我们回任务与联调记录看时间花在哪里，区分一项实现多花了多久，还是几端合起来以后补了检查。

\begin{table}[H]
  \centering
  \caption{核心任务三点估算（PERT）与实际工时偏差对照表}\label{tab:pert-estimation}
  \reporttablesize
  \begin{tabular}{@{}L{.10\linewidth}L{.255\linewidth}C{.07\linewidth}C{.07\linewidth}C{.07\linewidth}C{.09\linewidth}C{.08\linewidth}C{.08\linewidth}C{.09\linewidth}@{}}
\toprule
\textbf{任务} & \textbf{模块摘要} & \textbf{乐观$O$} & \textbf{可能$M$} & \textbf{悲观$P$} & \textbf{期望$E$(h)} & \textbf{标准差$\sigma$} & \textbf{实际(h)} & \textbf{偏差} \\
\midrule
T101--102 & SRS 建模与数据库基线 & 14 & 18 & 22 & 18.0 & 1.33 & 19.0 & +5.6\% \\
T103--105 & 认证、检索、购物车与下单事务 & 24 & 30 & 36 & 30.0 & 2.00 & 31.5 & +5.0\% \\
T106--107 & 顾客首页、详情、地址与结算页 & 18 & 24 & 30 & 24.0 & 2.00 & 24.5 & +2.1\% \\
T108 & Sprint 1 认证与下单测试用例 & 6 & 8 & 10 & 8.0 & 0.67 & 7.5 & -6.3\% \\
T201--203 & 商家履约、骑手抢单与资产服务 & 24 & 30 & 36 & 30.0 & 2.00 & 31.0 & +3.3\% \\
T204--205 & 商家端页面与骑手工作台 & 17 & 22 & 27 & 22.0 & 1.67 & 22.5 & +2.3\% \\
T206--207 & 备注/评价接口与集中联调缺陷单 & 12 & 16 & 20 & 16.0 & 1.33 & 16.5 & +3.1\% \\
T301--303 & 管理治理与多模态 AI 服务 & 22 & 28 & 34 & 28.0 & 2.00 & 29.0 & +3.6\% \\
T304 & 高德地图代理与坐标距离适配 & 8 & 10 & 12 & 10.0 & 0.67 & 10.5 & +5.0\% \\
T305--306 & AI 前端、主题动效与管理看板 & 19 & 24 & 29 & 24.0 & 1.67 & 25.0 & +4.2\% \\
T307--308 & TDD 导航修复与全链路交叉验收 & 15 & 20 & 25 & 20.0 & 1.67 & 21.0 & +5.0\% \\
\midrule
\textbf{合计} & \textbf{三大 Sprint 全周期汇总} & \textbf{179} & \textbf{230} & \textbf{281} & \textbf{230.0} & \textbf{---} & \textbf{238.0} & \textbf{+3.48\%} \\
\bottomrule
  \end{tabular}
\end{table}

\begin{figure}[!htbp]
  \centering
  \begin{minipage}[t]{.46\linewidth}
    \centering\vspace{0pt}
    \resizebox{.98\linewidth}{!}{%
\begin{tikzpicture}[x=1.6cm,y=0.075cm]
    \draw[->] (0,0) -- (5.2,0);
    \draw[->] (0,0) -- (0,78) node[above,font=\scriptsize] {工时 (h)};
    \foreach \y in {0,20,40,60}
      {\draw[gray!25] (0,\y) -- (4.8,\y); \node[left,font=\scriptsize] at (0,\y) {\y};}
    % 李承文
    \fill[tjublue] (0.6,0) rectangle (1.1,68); \node[above,font=\scriptsize] at (0.85,68) {68.0 h};
    \node[below,font=\scriptsize] at (0.85,0) {李承文};
    % 范文丽
    \fill[accentgreen!80!black] (1.6,0) rectangle (2.1,63); \node[above,font=\scriptsize] at (1.85,63) {63.0 h};
    \node[below,font=\scriptsize] at (1.85,0) {范文丽};
    % 赵紫怡
    \fill[accentorange] (2.6,0) rectangle (3.1,56); \node[above,font=\scriptsize] at (2.85,56) {56.0 h};
    \node[below,font=\scriptsize] at (2.85,0) {赵紫怡};
    % 仲禹潼
    \fill[tjublue!55] (3.6,0) rectangle (4.1,51); \node[above,font=\scriptsize] at (3.85,51) {51.0 h};
    \node[below,font=\scriptsize] at (3.85,0) {仲禹潼};
  \end{tikzpicture}}
    \captionof{figure}{四位成员全周期实际投入工时}\label{fig:mgmt-workload}
  \end{minipage}\hfill
  \begin{minipage}[t]{.51\linewidth}
    \centering\vspace{0pt}
    \captionof{table}{各轮成员实际投入工时（h）}\label{tab:member-workload}
    \reporttablesize
    \setlength{\tabcolsep}{3pt}
    \begin{tabular}{@{}lrrrr@{}}
      \toprule
      \textbf{成员} & \textbf{首轮} & \textbf{次轮} & \textbf{末轮} & \textbf{总计} \\
      \midrule
      李承文 & 23.0 & 22.0 & 23.0 & 68.0 \\
      范文丽 & 21.0 & 19.0 & 23.0 & 63.0 \\
      赵紫怡 & 21.5 & 15.0 & 19.5 & 56.0 \\
      仲禹潼 & 17.0 & 14.0 & 20.0 & 51.0 \\
      \midrule
      合计 & 82.5 & 70.0 & 85.5 & 238.0 \\
      \bottomrule
    \end{tabular}
  \end{minipage}
\end{figure}

\paragraph{各轮工时偏差}
Sprint 1、2、3 分别比计划多投入 2.5、2.0 和 3.5~h。首轮增加了历史菜名快照和相应回归，认证与下单测试则比估算少用 0.5~h。第二轮的新增工作主要是新人券判重和登录样式隔离。第三轮集中处理 AI 页返回循环和 375~px 窄屏溢出，并补充相应回归。表~\ref{tab:pert-estimation} 的差额按模块统计，同一次联调涉及的测试和修复也可从迭代记录回查。

表~\ref{tab:member-workload} 列出每轮的实际工作量，图~\ref{fig:mgmt-workload} 给出全周期对比。四人平均投入 59.5~h。李承文、范文丽、赵紫怡和仲禹潼分别占 28.6\%、26.5\%、23.5\% 和 21.4\%。第二轮主要集中在履约和资产联调，第三轮前端与测试的投入有所增加，与当轮看板、响应式统一、地图及交叉验收的安排相符。

\subsection{项目甘特图}

甘特图帮我们看哪几项可以同一天做，哪几项要等前一项。图~\ref{fig:mgmt-gantt} 按三轮的实际时间列了任务，四个里程碑的验收记录见表~\ref{tab:milestones}。

第一轮中，顾客端页面在认证和店铺接口开发期间已开始联调，订单快照与取消事务随后补入主链路。第二轮并行推进商家、骑手和资产服务，联调窗口落在这些任务后半段。图中横条的重叠能用来判断切入联调的时机，里程碑的计划日和验收日则按表~\ref{tab:milestones} 核对。

\begin{figure}[!htbp]
  \centering
  \resizebox{.94\linewidth}{!}{%
\begin{tikzpicture}[x=0.60cm,y=0.62cm]
    % 背景网格与日期刻度 (Day 1 = 9/1 ... Day 19 = 9/19)
    \foreach \d/\lbl in {1/9.1,3/9.3,5/9.5,7/9.7,9/9.9,11/9.11,13/9.13,15/9.15,17/9.17,19/9.19}
      {\draw[gray!25] (\d,0.5) -- (\d,12.5); \node[above,font=\scriptsize] at (\d,12.5) {\lbl};}
    % Sprint 背景分区
    \fill[blue!4] (0.5,0.5) rectangle (7.5,12.5);
    \fill[orange!4] (7.5,0.5) rectangle (12.5,12.5);
    \fill[green!4] (12.5,0.5) rectangle (19.5,12.5);
    \node[font=\scriptsize\bfseries,tjublue] at (4.0,14.8) {Sprint 1（核心闭环）};
    \node[font=\scriptsize\bfseries,accentorange] at (10.0,14.8) {Sprint 2（履约与资产）};
    \node[font=\scriptsize\bfseries,accentgreen] at (16.0,14.8) {Sprint 3（智能治理与验收）};

    % 任务条
    \node[left,font=\scriptsize] at (0.3,11.8) {SRS 建模与表结构};
    \fill[tjublue] (1,11.55) rectangle (3.5,12.05);
    \node[left,font=\scriptsize] at (0.3,10.8) {认证与店铺商品后端};
    \fill[tjublue] (2.5,10.55) rectangle (6,11.05);
    \node[left,font=\scriptsize] at (0.3,9.8) {订单快照与取消事务};
    \fill[tjublue] (4,9.55) rectangle (7.5,10.05);
    \node[left,font=\scriptsize] at (0.3,8.8) {顾客端核心页面联调};
    \fill[tjublue!70] (3.5,8.55) rectangle (7.5,9.05);

    \node[left,font=\scriptsize] at (0.3,7.8) {商家端履约与菜品管理};
    \fill[accentorange] (7.5,7.55) rectangle (11.5,8.05);
    \node[left,font=\scriptsize] at (0.3,6.8) {骑手抢单与配送状态机};
    \fill[accentorange] (8,6.55) rectangle (11.5,7.05);
    \node[left,font=\scriptsize] at (0.3,5.8) {钱包积分优惠券体系};
    \fill[accentorange!75] (8.5,5.55) rectangle (12,6.05);
    \node[left,font=\scriptsize] at (0.3,4.8) {集中联调与 WPS 缺陷单};
    \fill[accentorange!55] (10.5,4.55) rectangle (12.5,5.05);

    \node[left,font=\scriptsize] at (0.3,3.8) {管理端审核与统计看板};
    \fill[accentgreen] (12.5,3.55) rectangle (16.5,4.05);
    \node[left,font=\scriptsize] at (0.3,2.8) {多模态 AI 与高德地图};
    \fill[accentgreen!80] (13,2.55) rectangle (16.5,3.05);
    \node[left,font=\scriptsize] at (0.3,1.8) {TDD 修复与四套主题动效};
    \fill[accentgreen!65] (14.5,1.55) rectangle (17.5,2.05);
    \node[left,font=\scriptsize] at (0.3,0.8) {575 项回归与交叉验收};
    \fill[accentred!75] (16,0.55) rectangle (19.5,1.05);

    % 里程碑菱形
    \draw[fill=tjublue] (3.5,13.55) -- (3.72,13.80) -- (3.5,14.05) -- (3.28,13.80) -- cycle;
    \node[right,font=\scriptsize] at (3.75,13.80) {M1: SRS 2.0};
    \draw[fill=accentorange] (7.5,13.55) -- (7.72,13.80) -- (7.5,14.05) -- (7.28,13.80) -- cycle;
    \node[right,font=\scriptsize] at (7.75,13.80) {M2: MVP};
    \draw[fill=accentgreen] (12.5,13.55) -- (12.72,13.80) -- (12.5,14.05) -- (12.28,13.80) -- cycle;
    \node[right,font=\scriptsize] at (12.75,13.80) {M3: 联调封板};
    \draw[fill=accentred] (19.0,13.55) -- (19.22,13.80) -- (19.0,14.05) -- (18.78,13.80) -- cycle;
    \node[left,font=\scriptsize] at (18.75,13.80) {M4: 交付};
  \end{tikzpicture}
}
  \caption{轻量级外卖服务平台项目实施甘特图与里程碑节点}\label{fig:mgmt-gantt}
\end{figure}

\begin{longtable}{C{.10\linewidth}L{.22\linewidth}C{.13\linewidth}L{.46\linewidth}}
\caption{项目核心里程碑（M1--M4）检查与达成结果表}\label{tab:milestones}\\
\toprule
\textbf{里程碑} & \textbf{里程碑名称} & \textbf{计划/实际日} & \textbf{标志性准出交付物} \\
\midrule
M1 & SRS 2.0 需求冻结 & 9.6 / 9.6 & 17 张 UML/draw.io 建模图、23 张表 DDL 与 22 项用户故事基线。 \\
M2 & MVP 核心链路贯通 & 9.7 / 9.7 & 游客/顾客浏览、购物车、防重提交、取消回滚与登录鉴权跑通。 \\
M3 & 四端全链路联调封板 & 9.15 / 9.15 & 商家/骑手/管理端协同、资产积分体系、AI 多模态与统一壳层合入 \code{main}。 \\
M4 & 验收交付与文档归档 & 9.19 / 9.19 & 575 项全量回归通过、25 项交叉验收证据归档、Docker 一键演示就绪。 \\
\end{longtable}

M1 核定需求和结构，M2 核定可运行的顾客链路。两者间隔一天，因为文档、数据库与认证代码在第一轮并行推进。M2 之后，商家与骑手以已经建立的订单为起点联调。到 M3 时，团队把 AI、主题与四端页面集成在同一版本上，集中修复阻碍操作的问题。M4 再核对回归、交叉证据与演示环境。每个里程碑只检查该阶段应交付的结果，相关依据列在表~\ref{tab:milestones} 中。

\subsection{项目整体燃尽图}

22 项用户故事共 118 个故事点，其中 P0 为 78 点、P1 为 31 点、P2 为 9 点。我们在图~\ref{fig:mgmt-overall-burndown} 上按天记还剩多少点。前期补了 SRS 2.0 的模型和订单菜名快照，收敛稍慢。第二轮把认证、商品和履约主线连起来，9 月 15 日封板后继续验收，9 月 19 日清零。每轮的速率（Velocity）见表~\ref{tab:sprint-velocity}。

\begin{figure}[!htbp]
  \centering
  \resizebox{.86\linewidth}{!}{%
\begin{tikzpicture}[x=0.62cm,y=0.055cm]
    \draw[->] (0,0) -- (20.2,0) node[right,font=\scriptsize] {研发日 (9.1--9.19)};
    \draw[->] (0,0) -- (0,130) node[above,font=\scriptsize] {剩余故事点};
    \foreach \d/\lbl in {1/9.1,3/9.3,5/9.5,7/9.7,9/9.9,11/9.11,13/9.13,15/9.15,17/9.17,19/9.19}
      {\draw[gray!25] (\d,0) -- (\d,120); \node[below,font=\scriptsize] at (\d,0) {\lbl};}
    \foreach \y in {0,30,60,90,118}
      {\draw[gray!25] (0,\y) -- (19.5,\y); \node[left,font=\scriptsize] at (0,\y) {\y};}
    % 理想燃尽线
    \draw[dashed,mutedgray,thick] (1,118) -- (19,0);
    % 实际剩余故事点折线
    \draw[thick,tjublue,mark=*] plot coordinates {
      (1,118) (2,115) (3,108) (4,104) (5,96) (6,88) (7,76)
      (8,68) (9,60) (10,54) (11,45) (12,36) (13,30) (14,22)
      (15,12) (16,6) (17,3) (18,1) (19,0)
    };
    \node[font=\scriptsize,tjublue] at (13.5,95) {— 实际剩余故事点};
    \node[font=\scriptsize,mutedgray] at (13.5,108) {- - 理想燃尽基准线};
  \end{tikzpicture}
}
  \caption{项目全周期整体故事点燃尽图（总计 118 故事点，9 月 1 日--9 月 19 日）}\label{fig:mgmt-overall-burndown}
\end{figure}

到 9 月 7 日，图~\ref{fig:mgmt-overall-burndown} 剩余 76 点，对应首轮交付 42 点。9 月 12 日剩余 36 点，第二轮又交付 40 点。9 月 15 日封板时还剩 12 点，曲线在最后四天逐步收至零。把日线与表~\ref{tab:sprint-velocity} 的迭代汇总对照，可以区分每天的收敛节奏和每轮实际验收的总量。

\begin{table}[H]
  \centering
  \caption{三大 Sprint 敏捷速率（Velocity）与完成率统计表}\label{tab:sprint-velocity}
  \reporttablesize
  \begin{tabular}{@{}L{.18\linewidth}C{.16\linewidth}C{.16\linewidth}C{.16\linewidth}L{.25\linewidth}@{}}
\toprule
\textbf{迭代周期} & \textbf{计划故事点} & \textbf{实际交付点数} & \textbf{迭代完成率} & \textbf{主要交付增量} \\
\midrule
Sprint 1（7 天） & 42 点 & 42 点 & 100\% & SRS 2.0、23 张表基线、认证与下单 MVP \\
Sprint 2（5 天） & 40 点 & 40 点 & 100\% & 商家履约、骑手抢单、资产流水与评价 \\
Sprint 3（7 天） & 36 点 & 36 点 & 100\% & 管理治理、多模态 AI、高德适配与回归 \\
\midrule
\textbf{合计（19 天）} & \textbf{118 点} & \textbf{118 点} & \textbf{100\%} & \textbf{全部 P0/P1/P2 目标交付完毕} \\
\bottomrule
  \end{tabular}
\end{table}

我们把交付多少和用了多久分开看。第二轮 5 天验收了计划的 40 点，实际投入 70.0~h，比计划多 2.0~h。故事点燃尽与实际工时给出了这轮项目的进度和投入。要继续查为什么多花了时间，我们再看过程留下的联调问题、返工和回归记录，如新人券判重和登录样式隔离。这些过程数据有了指向，下一轮估算才能考虑真实发生的工作。

\FloatBarrier
\Needspace{6\baselineskip}
\section{风险评估与规避}

\subsection{风险识别与评估清单}

我们把容易出问题的地方与实际影响记进风险台账。比如模型响应慢会拖长一次调用，订单状态说法不一致则会让几端对同一笔订单给出不同结果。每项分别按高、中、低估计发生概率和影响，再排处置顺序。图~\ref{fig:risk-heatmap} 汇总了 12 项风险，预案见附录~\ref{app:process} 的表~\ref{tab:risk-register}。

\begin{figure}[H]
  \centering
  \resizebox{.72\linewidth}{!}{%
\begin{tikzpicture}[scale=1.0]
    % 3x3 风险热力网格
    \fill[yellow!20] (0,0) rectangle (2.6,1.2);
    \fill[orange!20] (2.6,0) rectangle (5.2,1.2);
    \fill[red!18]   (5.2,0) rectangle (7.8,1.2);
    \fill[green!12] (0,1.2) rectangle (2.6,2.4);
    \fill[yellow!22] (2.6,1.2) rectangle (5.2,2.4);
    \fill[orange!25] (5.2,1.2) rectangle (7.8,2.4);
    \fill[red!18]   (2.6,2.4) rectangle (5.2,3.6);
    \fill[red!30]   (5.2,2.4) rectangle (7.8,3.6);
    \fill[yellow!20] (0,2.4) rectangle (2.6,3.6);

    \draw[gray!60] (0,0) grid[xstep=2.6,ystep=1.2] (7.8,3.6);

    % 轴标签
    \node[left,font=\scriptsize\bfseries] at (-0.1,3.0) {概率：高};
    \node[left,font=\scriptsize\bfseries] at (-0.1,1.8) {概率：中};
    \node[left,font=\scriptsize\bfseries] at (-0.1,0.6) {概率：低};
    \node[below,font=\scriptsize\bfseries] at (1.3,-0.1) {影响：低};
    \node[below,font=\scriptsize\bfseries] at (3.9,-0.1) {影响：中};
    \node[below,font=\scriptsize\bfseries] at (6.5,-0.1) {影响：高};

    % 风险编号分布
    \node[align=center,font=\scriptsize\bfseries,accentred] at (6.5,3.0) {极高风险\\R001};
    \node[align=center,font=\scriptsize\bfseries,accentorange] at (3.9,3.0) {高风险\\R005, R009};
    \node[align=center,font=\scriptsize\bfseries,accentred] at (6.5,1.8) {高风险\\R002, R003, R004\\R007, R010, R011};
    \node[align=center,font=\scriptsize,tjublue] at (3.9,1.8) {中风险\\R006, R008, R012};
  \end{tikzpicture}
}
  \caption{项目 12 项关键风险概率--影响评估矩阵图}\label{fig:risk-heatmap}
\end{figure}

接口、状态含义不一致的 R001 落在高概率、高影响的一格。历史交易数据、库存并发与角色隔离等 6 项是中概率、高影响。模型调用和分支合并列在高概率、中影响一格，联调要给它们留回退或复核的时间。接下来按原编号在表~\ref{tab:risk-closure} 查处理结果和要继续观察的事项。

\subsection{风险跟踪与闭环}

我们把课上讲的风险缓解、监控和管理（RMMM）用在同一张台账里。能提前降低影响的，先定预案；正在发生变化的，留下检查结果；触发后怎么处理，也用原编号记录。交付时，逐项核对 12 项风险的工程处置或持续监控状态，见表~\ref{tab:risk-closure}。

\Needspace{.56\textheight}
\subsection{重点风险的处理与复查}

我们先处理会卡住多端联调的 R001。同一笔订单在顾客、商家和骑手端出现不同含义时，先确认 0--9 的状态对应，再把页面请求、后端 DTO 和契约测试放到一起看，新增接口也同时更新断言。R002、R003 和 R007 从数据上验，分别看历史快照有没有读原值、抢单与库存会不会被重复改动、配送费是否只在外送时计入。

模型调用的 R005 涉及超时和成本。我们以总时限和熔断限制等待，保留本地帮助，继续监控调用预算与限流。地图的 R006 直接用页面验。让地图服务不可用，再看顾客能否继续填写地址并提交订单。离线入口保留后，检查结果同处置状态一起写进风险记录。

交付前我们用风险原来描述的情形复查。权限从四端入口与账号隔离查起，部署在独立环境重启并访问健康端点，范围则回到本期 SRS。R010 提及的真实支付清算和实时 GPS 轨迹列入后续范围，本轮按模拟支付及已约定的履约流程验收。

\paragraph{同一风险的两层检查}
权限相关的 R004 和 R011 分别对应接口与页面。R011 检查各角色从哪一端登录，以及角色与目标路由不匹配时能否得到正确引导。R004 则在服务端确认资源归属，例如同为顾客的两个账号应读取各自的订单与购物车。验收时先走不同角色的页面，再针对同角色的不同账号验证数据。这样可以分清“入口跳转正确”和“后端读取范围正确”各由什么测试覆盖。

R008 的导航风险可以沿完整的返回路径复查。前端用例从原页面打开 AI 主入口，再进入工具页。第一次返回应回到 AI 主入口，第二次返回应回到原页面。用例也覆盖刷新后返回店铺列表，核对查询参数是否保留。页面展示另在窄屏检查，重点看结算按钮和输入框是否仍能操作。这样，修完返回循环后还可以用原用例防止它出现在其他入口。

对交付环境的 R012，代码合入时先检查本机绝对路径、配置文件与挂载路径。独立运行时再确认服务能够启动、健康端点正常，数据库及上传目录能使用准备好的环境配置。配置与运行两个结果分别留在仓库检查和部署验收中，复查时可以判断故障发生在文件打包、容器启动还是服务连接。

演示环境的首次启动可按数据库、迁移、应用的顺序排查。MySQL 就绪后，迁移容器应用订单名称快照和购物车备注两项迁移，成功后以退出码 0 结束。后端在迁移成功后启动，前端通过 Nginx 转发接口请求。首页已打开但订单加载失败时，可按这个顺序查看日志，区分代理转发、数据库连接和迁移故障。
